% Appendix: LLM coding instructions for the three Nairaland coding tasks.
% Each coder received the framing, the coding rules below, and one batch of
% posts (tid, page, k, post_year, title, text); coders returned one row per
% post in input order. Requires booktabs.

%% ---------------------------------------------------------------------------
\begin{table}[t]
\centering
\footnotesize
\begin{tabular}{p{2.2cm}p{11.6cm}}
\toprule
\multicolumn{2}{p{13.8cm}}{\textbf{Task framing.} ``You are coding Nairaland
(Nigerian forum) posts for stated family-size ideals. This is a full sweep:
most posts are ordinary conversation and state no ideal; expect the large
majority to be \texttt{has\_ideal} = 0. For each post, judge whether the
poster states their own ideal, desired, or recommended number of children, in
any phrasing: `I want 2 kids', `three is enough for any couple', `give me a
boy and a girl and I'm done', `4 kids max', `I go born 3 pikin', `no more
than 2', `stop at 4'.''} \\
\midrule
\textbf{Variable} & \textbf{Coding rule} \\
\midrule
\texttt{has\_ideal} & 1 only for a genuine own or recommended ideal number
(0 counts as a number if the poster genuinely wants no children). Counts of
existing children, other people's families, news, jokes, and questions
code 0. \\
\addlinespace
\texttt{ideal\_n} & The integer: childfree = 0; ``a boy and a girl'' = 2;
ranges take the floored midpoint; ``maximum of $N$'' = $N$. Empty if
\texttt{has\_ideal} = 0. \\
\addlinespace
\texttt{ideal\_kind} & \textit{own} (the poster's own ideal) or
\textit{prescriptive} (a number recommended for others). \\
\bottomrule
\end{tabular}
\caption{Coding instructions for stated ideal family size. Every post written
2010--2018 in the analysis corpus was coded individually; no keyword filter
was applied.}
\label{tab:codebook_ideals}
\end{table}

%% ---------------------------------------------------------------------------
\begin{table}[t]
\centering
\footnotesize
\begin{tabular}{p{3.4cm}p{10.4cm}}
\toprule
\multicolumn{2}{p{13.8cm}}{\textbf{Task framing.} ``You are coding the
reasoning behind stated family-size ideals in Nairaland posts. For each post,
code which rationales the post gives for its stated ideal. Categories are
multi-label; leave empty if the post states the ideal with no reasoning
(many are bare numeric answers, which is expected). Code only reasoning the
poster actually gives for their own stated ideal, not general topic content.
Formulaic phrases like `by God's grace' are not
\texttt{religion\_blessing}.''} \\
\midrule
\textbf{Category} & \textbf{Definition} \\
\midrule
\texttt{economic\_cost} & Affordability, school fees, poverty, ``this economy'' \\
\texttt{child\_quality} & Raise, train, or educate few children well \\
\texttt{career\_aspiration} & Poster's career, business, personal goals \\
\texttt{womens\_autonomy} & Women's independence, body, health, freedom \\
\texttt{family\_planning} & Contraception or spacing practicalities as reason \\
\texttt{religion\_blessing} & Children as God's gift or will, religious duty \\
\texttt{lineage\_continuity} & Family name, heirs, inheritance \\
\texttt{security\_status} & Children as old-age support, social standing, helpers \\
\texttt{modern\_western} & Comparisons to abroad, developed countries, modern times \\
\texttt{traditional\_norm} & Appeals to culture, tradition, how it is done \\
\bottomrule
\end{tabular}
\caption{Coding instructions for rationales. All posts stating an ideal were
coded; categories may overlap within a post.}
\label{tab:codebook_rationales}
\end{table}

%% ---------------------------------------------------------------------------
\begin{table}[t]
\centering
\footnotesize
\begin{tabular}{p{2.6cm}p{5.6cm}p{5.6cm}}
\toprule
\multicolumn{3}{p{13.8cm}}{\textbf{Task framing.} ``You are coding Nairaland
posts from family-planning and childbearing threads for cultural scripts
about fertility. Code three independent binary dimensions; a post can be 1 on
any, all, or none. Most posts are banter, news reactions, or unrelated
debate: code 0 on all dimensions for those; that is expected for the large
majority of rows.'' Endorsement and opposition were coded in two separate
passes over the same posts, with the rules below.} \\
\midrule
\textbf{Script} & \textbf{Endorsing (pass 1)} & \textbf{Opposing (pass 2)} \\
\midrule
Planned fertility & Frames fertility as something to be planned: child
spacing, limiting births, stopping at a number, timing births, family
planning endorsed as practice, waiting until ready. Merely mentioning the
topic or opposing planning = 0. & Rejects deliberate fertility control:
births should be left to God, nature, or fate; family planning dismissed as
wrong, foreign, or pointless. Neutral reporting or questions = 0. \\
\addlinespace
Modern contraception / medical authority & Invokes modern contraceptive
methods by name (pills, IUD, implants, injectables, condoms, sterilization,
emergency contraception) or expert or medical authority (doctors, nurses,
hospitals, WHO). Traditional or folk methods alone (herbs, withdrawal,
abstinence, prayer) = 0. & Rejects or warns against modern methods:
religious or moral objection (sin, promotes promiscuity), fear or claims of
harm (pills or IUDs cause infertility or cancer), or promoting natural or
folk methods explicitly as better than modern ones. Balanced side-effect
information while still endorsing use = 0. \\
\addlinespace
Female reproductive autonomy & Invokes female autonomy or intentionality in
reproductive decisions: a woman deciding or having the right to decide her
births or her body, spousal joint decision framed as her equal say, refusing
husband or family pressure. General women's empowerment talk without
reproductive-decision content = 0. & Asserts that reproductive decisions
belong to the husband or family, that a wife must submit or needs her
husband's permission, or condemns a woman for deciding her childbearing or
contraception herself. General sexism without reproductive-decision
content = 0. \\
\bottomrule
\end{tabular}
\caption{Coding instructions for fertility scripts. Every post written
2010--2018 in the family-planning and childbearing threads was coded on all
three dimensions, once for endorsement and once for opposition.}
\label{tab:codebook_scripts}
\end{table}
